\begin{textsample}{NEG Dim 3 – 2009\_c4 – Score 1.00 – 2009\_c4/enade\_2009\_com\_soc\_rp\_c9c58bdcc7.txt}  \label{ex:f3_neg_155}
Capítulo 4 Percepção sobre a Prova As percepções dos estudantes sobre a prova aplicada no ENADE/2009 na área de COMUNICAÇÃO SOCIAL – RELAÇÕES PÚBLICAS foram mensuradas por meio de nove itens que avaliaram desde o grau de dificuldade da prova até o tempo gasto para concluí -la.
As questões foram analisadas separando -se concluintes e ingressantes, e as percepções sobre a prova foram relacionadas com o desempenho dos estudantes e com a Região Geográfica de localização do curso6.
O desempenho dos estudantes dentro de cada Área/Subárea foi classificado em quartos ( 1/4 ) para os concluintes e para os ingressantes em separado ( ver figura para o esquema de representação gráfica destas grandezas ).
As notas finais foram colocadas em ordem ascendente e no primeiro quarto foram colocadas as 25 \% notas mais baixas, limitadas superiormente pelo percentil 25 que vem a ser o primeiro quartil.
No segundo quarto, as 25 \% seguintes, limitadas superiormente pelo percentil 50 que vem a ser a mediana.
No terceiro quarto, as 25 \% notas após as encontradas no segundo quarto, ou seja, as primeiras 25 \% acima da mediana, limitadas superiormente pelo percentil 75 que vem a ser o terceiro quartil.
No último quarto foram alocadas as 25 \% maiores notas.
Potencialmente, poderia existir uma relação entre a nota do estudante e a sua percepção sobre a prova.
Calculando as médias das respostas segundo os quartos de notas, seria possível reconhecer a existência ( ou não ) de algum gradiente nas respostas, por exemplo, uma melhor nota correspondendo a uma percepção menos negativa.
Este gradiente seria traduzido, por exemplo, no gráfico 4.1a, pela diminuição das áreas laranja e vermelha ( Difícil e Muito difícil ) com o aumento do desempenho ( indo do primeiro ao último quarto ). 6 Curso, neste relatório, refere -se à unidade de análise para o Conceito ENADE e é caracterizado pela combinação de IES, Área/Subárea ( quando pertinente ) e município de habilitação.
A medida tradicional de “ grau de discriminação ” comparou os resultados dos estudantes do primeiro quarto ( 25 \% piores estudantes ) aos do último quarto ( 25 \% melhores ).
Esta comparação, no gráfico disponibilizado, foi idêntica à comparação da primeira barra ( primeiro quarto ) com a última barra ( último quarto ) de cada grupo ingressante ou concluinte nos gráficos “ segundo quartos de nota final ” ( gráficos com números ímpares neste capítulo ).
Por outro lado, a disponibilização das informações segundo os quartos de nota permitiu verificar o comportamento de dois outros grupos intermediários que na análise somente dos extremos tem que ser suposto monotônico.
A seguir, serão apresentados os principais resultados válidos relativos aos nove itens avaliados. 4.1 Grau de dificuldade da prova 4.1.1 Formação Geral Os Gráficos 4.1a, 4.1b e 4.1c apresentam a distribuição das respostas à questão “ Qual o grau de dificuldade desta prova na parte de Formação Geral? ”, segundo quartos de nota final e grupo de estudantes concluintes e ingressantes que participaram do ENADE/2009 na área de COMUNICAÇÃO SOCIAL – RELAÇÕES PÚBLICAS.
No total de concluintes, 13,8 \% declararam que a prova na parte da Formação Geral foi Difícil/Muito difícil, contra um valor um pouco mais elevado para os ingressantes, 16,5 \%.
Tanto entre os concluintes, como entre os ingressantes, a proporção de indivíduos que classificaram esta parte como Difícil/Muito difícil diminui com o desempenho do estudante.
Por outro lado, a proporção que declarou a parte Fácil/Muito fácil foi um pouco maior entre os concluintes ( 23,3 \% ) do que entre os ingressantes ( 21,0 \% ).
Não se observou uma variação monotônica significativa desta proporção com o aumento no desempenho, nem entre os concluintes, nem entre os ingressantes.
Entre os concluintes, 63,0 \% classificaram esta parte com um nível Médio de dificuldade.
A proporção entre os ingressantes foi ligeiramente menor, 62,5 \%.
Os Gráficos 4.2a e 4.2b apresentam a distribuição das respostas à questão “ Qual o grau de dificuldade desta prova na parte de Formação Geral? ”, segundo a Região onde o curso está situado os estudantes concluintes e ingressantes que participaram do ENADE/2009 na área de COMUNICAÇÃO SOCIAL – RELAÇÕES PÚBLICAS.
A Região Nordeste apresentou o menor percentual ( 7,8 \% ) de declaração de Difícil/Muito difícil.
Entre os concluintes, a maior proporção de estudantes que considerou a parte de Formação Geral da prova como de dificuldade média entre todas as Regiões brasileiras, foram os da Região Norte com 67,3 \%.
Entre os ingressantes foram do Nordeste, com 69,5 \%.
As Regiões com as menores proporções desta resposta foram para os concluintes a Sudeste ( 61,3 \% ) e a Sul ( 63,3 \% ) e para os ingressantes, as Regiões com as menores proporções desta resposta foram a Sul ( 60,9 \% ) e a Centro-Oeste ( 60,7 \% ). 4.1.2 Componente Específico Os Gráficos 4.3a, 4.3b e 4.3c apresentam a distribuição das respostas à questão “ Qual o grau de dificuldade desta prova na parte do Componente Específico? ”, segundo quartos de nota final e grupo de estudantes concluintes e ingressantes que participaram do ENADE/2009 na área de COMUNICAÇÃO SOCIAL – RELAÇÕES PÚBLICAS.
Em geral, para as desagregações consideraram, os estudantes que classificaram o Componente Específico como mais difícil que a Formação Geral.
No total de concluintes, 11,9 \% declararam que a prova na parte do Componente Específico foi Difícil/Muito difícil, sendo um valor muito maior para os ingressantes, 25,0 \%.
A proporção que declarou Difícil/Muito difícil diminui com o desempenho entre os concluintes e ingressantes.
A proporção que declarou a parte Fácil/Muito fácil foi maior entre os concluintes ( 25,2 \% ) do que entre os ingressantes ( 13,4 \% ).
Entre os concluintes, 63,0 \% classificaram esta parte com um nível Médio de dificuldade.
A proporção entre os ingressantes foi menor, 61,7 \%.
Entre os ingressantes esta proporção aumentou com o desempenho.
Os Gráficos 4.4a e 4.4b apresentam a distribuição das respostas à questão “ Qual o grau de dificuldade desta prova na parte do Componente Específico? ”, segundo a Região onde o curso está situado para os estudantes concluintes e ingressantes que participaram do ENADE/2009 na área de COMUNICAÇÃO SOCIAL – RELAÇÕES PÚBLICAS.
Entre os concluintes e ingressantes, a maior proporção de estudantes que considerou a parte do Componente Específico da prova como de dificuldade média entre todas as Regiões brasileiras, foi na Região Norte para os concluintes ( 73,1 \% ) e na Região Centro-Oeste para os ingressantes ( 70,5 \% ).
A Região com as menores proporções desta resposta foram respectivamente para os concluintes e ingressantes, a Sudeste ( 61,4 \% ) e a Sul ( 57,8 \% ). 4.2 Extensão da prova em relação ao tempo total Os Gráficos 4.5a, 4.5b e 4.5c apresentam a distribuição das respostas à questão sobre a extensão da prova em relação ao tempo total segundo quartos de nota final e grupo de estudantes concluintes e ingressantes que participaram do ENADE/2009 na área de COMUNICAÇÃO SOCIAL – RELAÇÕES PÚBLICAS.
Tanto para concluintes quanto para ingressantes, a maioria considerou a extensão da prova como Adequada em relação ao tempo destinado à sua resolução, respectivamente 60,7 \% e 62,0 \%.
Entre os concluintes e ingressantes esta proporção cresceu com o desempenho.
A proporção dos que consideraram a extensão da prova Muito curta foi ligeiramente maior entre os concluintes ( 0,9 \% ) do que entre os ingressantes ( 0,8 \% ).
A proporção dos que consideraram a extensão da prova Curta/Muito curta em relação ao tempo total foi menor entre os concluintes ( 4,8 \% ) do que entre os ingressantes ( 5,5 \% ).
Entre os concluintes, os estudantes com desempenho no grupo mais baixo tiveram o maior percentual na declaração de que a extensão da prova teria sido Muito curta com relação ao tempo total.
Entre os ingressantes este comportamento se mantém.
No entanto, entre os estudantes que consideraram a prova Curta, tanto para concluintes quanto para ingressantes, o maior percentual está entre aqueles que obtiveram os melhores desempenhos.
A proporção dos que consideraram a extensão da prova Muito longa foi entre os concluintes ( 12,7 \% ) e ingressantes ( 10,3 \% ).
A proporção dos que consideraram a extensão da prova Longa foi menor entre os concluintes ( 21,8 \% ) e os ingressantes ( 22,1 \% ).
Os Gráficos 4.6a e 4.6b apresentam a distribuição das respostas à questão sobre a extensão da prova em relação ao tempo total, segundo a Região onde o curso está situado para os estudantes concluintes e ingressantes que participaram do ENADE/2009 na área de COMUNICAÇÃO SOCIAL – RELAÇÕES PÚBLICAS.
Em todas as Regiões brasileiras, a maioria dos estudantes, tanto concluintes como ingressantes, perceberam a prova como Os Gráficos 4.6a e 4.6b apresentam a distribuição das respostas à questão sobre a extensão da prova em relação ao tempo total, segundo a Região onde o curso está situado para os estudantes concluintes e ingressantes que participaram do ENADE/2009 na área de COMUNICAÇÃO SOCIAL – RELAÇÕES PÚBLICAS.
Em todas as Regiões brasileiras, a maioria dos estudantes, tanto concluintes como ingressantes, perceberam a prova como Adequada.
Somente os concluintes das Regiões Norte e Centro-Oeste ( 63,9 \% e 69,5 \%, respectivamente ) se situaram acima da média ( 60,7 \% ) para esta resposta, enquanto entre os ingressantes, os da Região Norte ( 73,8 \% ) e Centro-Oeste ( 69,4 \% ) ficaram acima da média ( 62,0 \% ). 4.3 Compreensão dos Enunciados das Questões 4.3.1 Formação Geral Os Gráficos 4.7a, 4.7b e 4.7c apresentam a distribuição das respostas à questão “ Os enunciados das questões da prova na parte de Formação Geral, estavam claros e objetivos? ”, segundo quartos de nota final e grupo de estudantes concluintes e ingressantes que participaram do ENADE/2009 na área de COMUNICAÇÃO SOCIAL – RELAÇÕES PÚBLICAS.
Com relação aos enunciados das questões da parte de Formação Geral, as opiniões foram positivas : a grande maioria dos concluintes ( 76,5 \% ) e ingressantes ( 80,0 \% ) consideram claros e objetivos os enunciados de todas ou da maioria das questões.
Entre os concluintes, esta proporção cresceu com o desempenho : 64,0 \% para o primeiro grupo e 82,3 \% para o último grupo.
Entre os ingressantes o comportamento se repetiu o primeiro grupo, o de pior desempenho com 68,3 \% e o último grupo, o de melhor desempenho, com 89,1 \%.
Em linhas gerais, pode -se dizer que a combinação de poucos e nenhum decresceu com o aumento do desempenho tanto para concluintes como para ingressantes.
Os Gráficos 4.8a e 4.8b apresentam a distribuição das respostas à questão “ Os enunciados das questões da prova na parte de Formação Geral, estavam claros e objetivos? ”, segundo a Região onde o curso está situado para os estudantes concluintes e ingressantes que participaram do ENADE/2009 na área de COMUNICAÇÃO SOCIAL – RELAÇÕES PÚBLICAS.
Em termos regionais, as proporções de ingressantes e de concluintes que consideraram claros e objetivos os enunciados de todas ou da maioria das questões foram respectivamente, 80,0 \% e 76,4 \%.
Entre os concluintes, os valores foram de 68,4 \% no Centro-Oeste a 78,4 \% no Nordeste.
Já entre os ingressantes, estas proporções variaram de 77,3 \% no Sul a 87,8 \% no Norte.
Entre os ingressantes as maiores proporções ocorreram nas Regiões Nordeste e Sul, respectivamente 11,5 \% e 11,2 \%. 4.3.2 Componente Específico Os Gráficos 4.9a, 4.9b e 4.9c apresentam a distribuição das respostas à questão “ Os enunciados das questões da prova na parte de formação específica estavam claros e objetivos? ”, segundo quartos de nota final e grupo de estudantes concluintes e ingressantes que participaram do ENADE/2009 na área de COMUNICAÇÃO SOCIAL – RELAÇÕES PÚBLICAS.
Com relação aos enunciados das questões da parte do Componente Específico, as opiniões foram também positivas : a grande maioria dos concluintes ( 78,3 \% ) e ingressantes ( 80,2 \% ) consideraram claros e objetivos os enunciados de todas ou da maioria das questões.
O número de estudantes que declararam que “ Não, nenhum ” dos enunciados das questões da parte do Componente Específico era claro e objetivo, foi semelhante à parte de Formação Geral ( 1,5 \% entre os concluintes e 1,2 \% entre os ingressantes ) : 1,9 \% para concluintes e 1,2 \% para ingressantes.
Os que responderam “ Poucos ” são 7,4 \% entre os concluintes e 5,9 \% entre os ingressantes.
A proporção de estudantes que respondeu que “ cerca de metade ” dos enunciados da parte do Componente Específico era claro e objetivo decresceu com o desempenho, para ingressantes.
Os Gráficos 4.10a e 4.10b apresentam a distribuição das respostas à questão “ Os enunciados das questões da prova na parte de formação específica estavam claros e objetivos? ”, segundo a Região onde o curso está situado para os estudantes concluintes e ingressantes que participaram do ENADE/2009 na área de COMUNICAÇÃO SOCIAL – RELAÇÕES PÚBLICAS.
Em termos regionais, as proporções de ingressantes e de concluintes que consideraram claros e objetivos os enunciados de todas ou da maioria das questões variaram mais entre os ingressantes ( no máximo 8 pontos percentuais ) do que entre os concluintes ( no máximo 4 pontos percentuais ).
Entre os concluintes, os valores foram de 76,9 \% no Sudeste a 81,2 \% no Nordeste. “ Para os ingressantes, as Regiões Norte ( 7,1 \% ), Sudeste ( 6,4 \% ) e Centro-Oeste ( 4,8 \% ) apresentaram valores abaixo da média ( 7,2 \% ). ” 4.4 Suficiência das informações/instruções fornecidas Os Gráficos 4.11a, 4.11b e 4.11c apresentam a distribuição das respostas à questão “ As informações/instruções fornecidas para a resolução das questões foram suficientes para resolvê -las? ”, segundo quartos de nota final e grupo de estudantes concluintes e ingressantes que participaram do ENADE/2009 na área de COMUNICAÇÃO SOCIAL – RELAÇÕES PÚBLICAS.
Quando perguntados sobre as informações/instruções fornecidas para a resolução das questões, 85,8 \% dos ingressantes e 86,4 \% dos concluintes de todo o Brasil responderam que estas eram ou excessivas ou suficientes em todas ou na maioria das questões, demonstrando uma avaliação positiva da prova.
A proporção de estudantes com a opinião supracitada foi crescente com o desempenho, tanto para concluintes como para ingressantes.
Entre os concluintes a proporção vai de 76,6 \% para os de pior desempenho a 91,3 \% para os de melhor desempenho.
Entre os ingressantes, as proporções equivalentes foram, respectivamente, 76,1 \% a 94,2 \%.
O complemento desta opinião, ou seja, os estudantes que responderam “ Somente em algumas ” ou “ Em nenhuma delas ”, apresentou uma proporção decrescente com o desempenho tanto para os concluintes como para os ingressantes.
Os Gráficos 4.12a e 4.12b apresentam a distribuição das respostas à questão “ As informações/instruções fornecidas para a resolução das questões foram suficientes para resolvê -las? ”, segundo a Região onde o curso está situado para os estudantes concluintes e ingressantes que participaram do ENADE/2009 na área de COMUNICAÇÃO SOCIAL – RELAÇÕES PÚBLICAS.
Os Gráficos 4.12a e 4.12b apresentam a distribuição das respostas à questão “ As informações/instruções fornecidas para a resolução das questões foram suficientes para resolvê -las? ”, segundo a Região onde o curso está situado para os estudantes concluintes e ingressantes que participaram do ENADE/2009 na área de COMUNICAÇÃO SOCIAL – RELAÇÕES PÚBLICAS. 4.5 Dificuldade encontrada ao responder à prova Os Gráficos 4.13a, 4.13b e 4.13c apresentam a distribuição das respostas à questão “ Você se deparou com alguma dificuldade ao responder à prova? ”, segundo quartos de nota final e grupo de estudantes concluintes e ingressantes que participaram do ENADE/2009 na área de COMUNICAÇÃO SOCIAL – RELAÇÕES PÚBLICAS.
Perguntados sobre as dificuldades com as quais se depararam ao responder à prova, 30,8 \% dos ingressantes apontaram como principal dificuldade o desconhecimento do conteúdo, enquanto que os concluintes foram a forma diferente de abordagem, 31,5 \%.
Na verdade, para a quase totalidade das áreas, ingressantes apresentaram uma proporção de estudantes bem maior do que entre os concluintes declarando como principal dificuldade o desconhecimento da matéria.
O desconhecimento do conteúdo foi apontado somente por 8,1 \% dos concluintes, valores variando de 12,7 \% ( pior desempenho ) a 6,2 \% ( melhor desempenho ).
A forma diferente de abordagem do conteúdo foi considerada como a segunda maior dificuldade por 27,4 \% dos ingressantes, enquanto que a falta de motivação foi a segunda maior dificuldade entre os concluintes ( 30,8 \% ) sendo que esta opção foi apontada por 19,8 \% dos ingressantes.
Esta dificuldade apresentou valores no entorno de 34 \% para os três primeiros grupos de desempenho de concluintes.
Em termos nacionais, a proporção de ingressantes e concluintes que declararam não ter tido qualquer dificuldade para responder à prova foi, respectivamente, 14,7 \% e 19,7 \%.
Os Gráficos 4.14a e 4.14b apresentam a distribuição das respostas à questão “ Você se deparou com alguma dificuldade ao responder à prova? ”, segundo a Região de onde o curso está situado para os estudantes concluintes e ingressantes que participaram do ENADE/2009 na área de COMUNICAÇÃO SOCIAL – RELAÇÕES PÚBLICAS.
Nas Regiões à prova vinha do desconhecimento do conteúdo, sendo que na Região Norte, as maiores dificuldades para responder à prova vinham do desconhecimento do conteúdo e da forma diferente de abordagem do conteúdo, em proporções iguais ( 35,7 \% ).
Na Região Centro-Oeste foi a forma diferente de abordagem do conteúdo.
Entre os concluintes, a resposta à forma diferente de abordar o conteúdo variou de 26,7 \% na Região Nordeste a 35,9 \% na Região Sul.
Entre os ingressantes, essa opção variou de 25,8 \% na Região Centro-Oeste a 35,7 \% na Norte.
O motivo que teve o menor percentual de escolha foi o Desconhecimento do Conteúdo apontado por apenas 8,1 \% dos concluintes, tendo pouca variação entre as Regiões ; pois o menor percentual foi de 7,4 \% ( Região Centro-Oeste ) e o maior percentual foi de 9,7 \% ( Região Sul ).
O motivo que teve o menor percentual de escolha foi o Espaço insuficiente para responder às questões apontadas por 7,3 \% dos ingressantes variando de 4,8 \% ( Região Norte ) e o maior percentual foi de 11,3 \% ( Região Centro-Oeste ). 4.6 Percepção sobre os conteúdos das questões objetivas da prova Os Gráficos 4.15a, 4.15b e 4.15c apresentam a distribuição das respostas à questão “ Considerando apenas as questões objetivas da prova, você percebeu que :... ” segundo quartos de nota final e grupo de estudantes concluintes e ingressantes que participaram do ENADE/2009 na área de COMUNICAÇÃO SOCIAL – RELAÇÕES PÚBLICAS.
As respostas possíveis para esta questão estavam relacionadas à cobertura do conteúdo das questões objetivas durante o curso.
A proporção de ingressantes ( 41,5 \% ) que declarou que não estudou ainda a maioria desses conteúdos foi quase 10 vezes maior do que entre os concluintes ( 4,3 \% ).
Esta proporção foi decrescente tanto para os concluintes como para os ingressantes como função do desempenho.
Por outro lado, a proporção de estudantes que declarou que estudou e aprendeu todos ou muitos desses conteúdos, entre os ingressantes, foi de 40,4 \% e entre os concluintes 80,1 \%.
Esta proporção foi crescente como função do desempenho, tanto para os concluintes como para os ingressantes.
A proporção de estudantes que declarou ter estudado a maioria desses conteúdos, mas não os ter aprendido é maior entre os concluintes ( 8,6 \% ) do que os ingressantes ( 5,8 \% ).
Os Gráficos 4.16a e 4.16b apresentam a distribuição das respostas à questão “ Considerando apenas as questões objetivas da prova, você percebeu que :... ” segundo a Região onde o curso para os estudantes concluintes e ingressantes que participaram do ENADE/2009 na área de COMUNICAÇÃO SOCIAL – RELAÇÕES PÚBLICAS.
Em todas as Regiões, a maioria dos concluintes afirmou ter estudado e aprendido todos ou muitos dos conteúdos, com percentuais que variaram de 72,9 \% na Região Norte a 84,9 \%, na Centro-Oeste.
Entre os ingressantes, o comportamento foi semelhante com as porcentagens variando de 26,8 \% na Região Norte a 57,4 \% na Centro-Oeste. 4.7 Tempo gasto para concluir a prova Os Gráficos 4.17a, 4.17b e 4.17c apresentam a distribuição das respostas à questão “ Qual foi o tempo gasto para concluir a prova? ”, segundo quartos de nota final e grupo de estudantes concluintes e ingressantes que participaram do ENADE/2009 na área de COMUNICAÇÃO SOCIAL – RELAÇÕES PÚBLICAS.
Quando perguntados quanto ao tempo de conclusão da prova, a maioria respondeu entre uma e duas horas, tanto entre os concluintes ( 49,3 \% ) quanto entre os ingressantes ( 47,8 \% ).
A segunda maior incidência de respostas foi entre duas e três horas, igualmente para ambos os grupos : concluintes ( 30,5 \% ) e ingressantes ( 32,5 \% ).
A proporção daqueles que declararam ter usado as quatro horas, mas não terem conseguido terminar foi bem baixa : 0,8 \% entre os ingressantes e 1,4 \% entre os concluintes.
Entre os concluintes e ingressantes esta resposta ficou concentrada foi maior entre os estudantes do primeiro quarto, aqueles com pior desempenho.
Os Gráficos 4.18a e 4.18b apresentam a distribuição das respostas à questão “ Qual foi o tempo gasto para concluir a prova? ”, segundo a Região onde o curso está situado para os estudantes concluintes e ingressantes que participaram do ENADE/2009 na área de COMUNICAÇÃO SOCIAL – RELAÇÕES PÚBLICAS.
A proporção de concluintes que gastaram entre uma e duas horas foi maior na Região Norte ( 53,7 \% ) e menor na Região Sudeste ( 48,7 \% ).
Para ingressantes, o maior percentual para este tempo ocorreu na Região Nordeste ( 54,6 \% ) e a menor na Região Norte ( 43,9 \% ).

% matched lemmas: 
\end{textsample}
